../cictikz/src/cictikz/data/tex/cictikz_preamble.tex

% The non-overlap clock generator, and the waveforms it produces.
%
% Two cross-coupled NOR gates. Each output is fed back through a delay of two
% inverters into the other gate, so neither output can go high until the other
% has gone low and the delay has run out. That gap is the whole point of the
% circuit, so the right hand panel marks it.
%
% The original draws the two phases in green and red. They sit on separate
% baselines and carry their own labels, so the colour is not doing any work and
% both are drawn in black here.
%
% The delay chain is two inverters, not one: an odd number would invert the
% feedback and latch the generator.

\newcommand{\novlT}{1.9}         % clock period in the waveform panel
\newcommand{\novlX}{12.2}        % left edge of the waveform panel
\newcommand{\novlH}{1.0}         % logic swing

% One square wave: baseline, rising edge as a fraction of the period, falling
% edge as a fraction of the period. Three periods.
\newcommand{\novlWave}[3]{%
  \foreach \k in {0,1,2} {
    \draw[thick] ({\novlX+\k*\novlT},#1) -- ({\novlX+(\k+#2)*\novlT},#1)
      -- ({\novlX+(\k+#2)*\novlT},#1+\novlH)
      -- ({\novlX+(\k+#3)*\novlT},#1+\novlH)
      -- ({\novlX+(\k+#3)*\novlT},#1)
      -- ({\novlX+(\k+1)*\novlT},#1);
  }
}

% A labelled terminal.
\newcommand{\novlPort}[3]{%
  \draw[thick] (#1) circle (0.09);
  \node[anchor=#2] at (#1) {#3};
}

\begin{circuitikz}[american, thick, transform shape, circuit ee IEC]

  ../cictikz/src/cictikz/data/tex/cictikz_lib.tex

  % ---- The two NOR gates, house style ----
  \draw (4.2,2.4) \cicNor{n1};
  \draw (4.2,-2.4) \cicNor{n2};

  % ---- Clock in, straight to the top gate and inverted to the bottom ----
  \draw (0.3,2.72) circle (0.09);
  \node[anchor=east] at (0.15,2.72) {$\mathrm{Clk}$};
  \draw (0.39,2.72) -- (n1_in1);
  \fill (1.3,2.72) circle (0.07);
  \begin{scope}[shift={(1.3,2.1)}, rotate=-90]
    \draw (0,0) \cicInv{ck};
  \end{scope}
  \draw (1.3,2.72) -- (ck_in);
  \draw (ck_out) -- (1.3,-2.72) -- (n2_in2);

  % ---- Outputs ----
  \draw (n1_out) -- (10.4,2.4);
  \draw (10.49,2.4) circle (0.09);
  \node[anchor=west] at (10.65,2.4) {$\phi_{1}$};
  \draw (n2_out) -- (10.4,-2.4);
  \draw (10.49,-2.4) circle (0.09);
  \node[anchor=west] at (10.65,-2.4) {$\phi_{2}$};

  % ---- phi_1 delayed by two inverters into the bottom gate ----
  \fill (9.4,2.4) circle (0.07);
  \draw (9.4,2.4) -- (9.4,1.0);
  \draw (9.4,1.0) \cicInvM{d1a};
  \draw (d1a_out) -- (7.9,1.0) \cicInvM{d1b};
  \draw (d1b_out) -- (3.6,1.0) -- (3.6,-2.08) -- (n2_in1);

  % ---- phi_2 delayed the same way into the top gate ----
  \fill (9.4,-2.4) circle (0.07);
  \draw (9.4,-2.4) -- (9.4,-1.0);
  \draw (9.4,-1.0) \cicInvM{d2a};
  \draw (d2a_out) -- (7.9,-1.0) \cicInvM{d2b};
  \draw (d2b_out) -- (3.2,-1.0) -- (3.2,2.08) -- (n1_in2);

  % ---- Waveforms ----
  % phi_1 high over [0.05T, 0.45T], phi_2 over [0.55T, 0.95T]: a tenth of a
  % period of dead time at every handover, in both directions.
  \novlWave{1.5}{0.05}{0.45}
  \node[anchor=east] at (\novlX-0.15,1.5+0.5*\novlH) {$\phi_{1}$};
  \novlWave{-1.5}{0.55}{0.95}
  \node[anchor=east] at (\novlX-0.15,-1.5+0.5*\novlH) {$\phi_{2}$};

  % The gap between phi_1 falling and phi_2 rising.
  \draw[dotted, thick] ({\novlX+1.45*\novlT},-2.0) -- ({\novlX+1.45*\novlT},3.4);
  \draw[dotted, thick] ({\novlX+1.55*\novlT},-2.0) -- ({\novlX+1.55*\novlT},3.4);
  \draw[thick, <-] ({\novlX+1.5*\novlT},3.5) -- ++(0.5,0.6);
  \node[anchor=west] at ({\novlX+1.5*\novlT+0.55},4.15) {Non-overlap};

\end{circuitikz}
\end{document}
